\hypertarget{what-the-collective-circuit-costs-to-compile}{%
\subsection{What the collective circuit costs to compile}\label{sec:compilation}}

The collective route's overhead is often quoted as the reason not to use it, so it is worth
separating what is a compilation fact from what is a device measurement. The counts below are the
former: they are what a transpiler returns for a stated circuit against a published coupling map,
and they need no device access to reproduce.

For \(k = 2\) the destructive test adds \(n\) two-qubit gates and depth two on top of state
preparation, with no ancilla. Transpiled onto a published square-lattice map with a native
\(\{\mathrm{CZ}, R_x, R_z\}\) basis, the \(n = 2\) Bell instance uses four CZ gates --- two to
prepare each copy and two for the pairing --- with zero routing overhead, and the GHZ ladders at
\(n = 3, 4\) map with zero routing at \(3n - 2\) CZ gates and nest inside one another. Structured
states on a matched topology are therefore cheap to pair; generic states are not, and how much
they cost depends on the state-preparation synthesis, so we make no claim there.

For \(k \ge 3\) the ancilla-based Hadamard test is the relevant circuit and it is not comparable.
Transpiled to the same map, the controlled cyclic shift on \(kn\) qubits uses about ten times the
two-qubit gate count of the \(k = 2\) destructive test at \(k = 3\) and about sixteen times at
\(k = 4\), and needs an ancilla and routing the destructive test does not. Every \(k \ge 3\) result
in this paper is a simulation of an idealized bounded-outcome primitive that carries none of that
cost, and the reader should read those sizes as an upper bound on what the collective route buys.

These are gate counts, not error budgets. What they establish is that the entangling overhead is
not the binding constraint for the structured \(k = 2\) circuits we study; what binds instead is
the statistical failure of the single-copy estimator, which Section~\ref{sec:where-useless}
quantifies. A companion study takes up how the collective route behaves on a real device, where
neither gate counts nor an idealized primitive are the whole story.
